\scanpage{63}
On travaillera dorénavant dans la catégorie $P_1=Q$ (équivalente à la catégorie $Q_{\bullet}$ des complexes de chaînes finis où les \deleted{\unsure{-}} $L_i$, $Z_i$, $B_i$, $H_i$ sont libres de type fini), \inserted{munie de $\Sigma_1=\Sigma\cap P_1$,} dont la sous-catégorie des $\Sigma$-projectifs est \deleted{ad} la catégorie \inserted{additive} $P_0$ engendrée par les $\mathbb{C}^i$ et $\Psi^0$ (i.e. par les $\Phi^i$), et celle des $\Sigma$-injectifs est la sous-cat. additive $\mathbb{K}$ \deleted{\unsure{des faisceaux}} \inserted{engendrée} par les $\mathbb{C}^i$ (les $L_{\bullet}$ \deleted{flasques} acycliques i.e. les $L_*$ flasques --). \uline{On s'intéressera aux résolutions droites} \inserted{$K^{\bullet}$} $\Sigma$-\uline{injectives de} $\Psi^0$, (i.e. des résolutions flasques $K^{\bullet}_*$ de $k_*$ \unsure{$\simeq$} -- ce qui donneront, \inserted{du} en top algébrique \inserted{\unsure{$\simeq$}}, des méthodes de calcul remarquables de la coh à coeffs dans $k$). \unsure{On a|On a.} d'algèbre homologique ``relative'' \unsure{au choix de|au choix des|qui donne des} suites exactes, qui nous disent que si (comme c'est le cas ici) tout objet \deleted{\unsure{est q}} se plonge dans un $\Sigma$-injectif, \deleted{de façon} par un $\Sigma$-mono, alors \deleted{les} \deleted{les deux résolutions sont canoniquement \unsure{iso.}} deux résolutions sont homotopes, et il y a une seule classe d'homotopie d'hom de l'une dans l'autre (\uline{NB} Il y a aussi assez de projectifs, et on a l'assertion analogue pour les résolutions projectives; on fera attention que la cat. $\mathbb{K}$ des sommes de $\mathbb{C}^i$ \unsure{sera|bien|est}
% archive sheet number: 62
